\documentclass[12pt,a4paper]{report}
% NON MODIFICARE il formato del documento: è quello standard per il frontespizio.
\usepackage[italian]{babel}
\usepackage{newlfont}
\usepackage{graphicx}
\usepackage[top=2cm]{geometry}
\textwidth=450pt\oddsidemargin=0pt

\begin{document}
\begin{titlepage}
\begin{center}

% ------------------------------------------------------------------
% LOGO DELL'UNIVERSITÀ
% ------------------------------------------------------------------
% Prima di compilare:
% 1. Scarica il logo ufficiale dell'Università di Bologna dal sito di Ateneo.
% 2. Salva l'immagine nella stessa cartella del file .tex con il nome:
%       logo_unibo.png
%    (oppure modifica il nome qui sotto per farlo coincidere con il file che hai scaricato).
% 3. Assicurati che l'immagine sia in formato .png, .jpg o .pdf.
\includegraphics[width=0.25\linewidth]{logo_unibo.png}\vspace{0.7cm}\\

% Nome ufficiale dell'Ateneo: NON modificare.
{{\Large{\textsc{Alma Mater Studiorum $\cdot$ Universit\`a di
Bologna}}}} \rule[0.1cm]{15.8cm}{0.1mm}
\rule[0.5cm]{15.8cm}{0.6mm}

% ------------------------------------------------------------------
% DIPARTIMENTO / CORSO DI LAUREA
% ------------------------------------------------------------------
% Selezionare il corso di Laurea di riferimento, cancellando le opzioni non necessarie.
%
{\small{\bf Dipartimento di Informatica --- Scienza e Ingegneria\\
Corso di Laurea Magistrale/Triennale in Informatica/Informatica per il Management}}
\end{center}

\vspace{15mm}
\begin{center}
% ------------------------------------------------------------------
% TITOLO DELLA TESI
% ------------------------------------------------------------------
% Sostituire "TITOLO DELLA TESI" con il titolo completo della propria tesi.
% Suggerimenti:
% - Usare MAIUSCOLO per il titolo principale, se richiesto dalle linee guida.
% - Evitare abbreviazioni non standard.
{\LARGE{\bf TITOLO DELLA TESI}}\\
\vspace{3mm}

\vspace{19mm}
% ------------------------------------------------------------------
% SETTORE SCIENTIFICO DI RIFERIMENTO (OPZIONALE)
% ------------------------------------------------------------------
% Questa riga è opzionale. Puoi:
% - Sostituire "Settore Scientifico di Riferimento" con il nome dell'insegnamento
%   o ambito disciplinare da cui nasce la tesi (es. "Machine Learning", "Calcolo Numerico").
% - Oppure, se non richiesto dal tuo Corso di Laurea, puoi COMMENTARE
%   l'intera riga anteponendo un % all'inizio.
{\large{\bf Tesi di Laurea in ''Settore Scientifico di Riferimento'' \textit{(Opzionale)}}}
\end{center}

\vspace{30mm}
\par
\noindent
\begin{minipage}[t]{0.50\textwidth}
% ------------------------------------------------------------------
% RELATORE
% ------------------------------------------------------------------
% Sostituire:
% - "Chiar.ma Prof." con il titolo corretto del relatore:
%    * "Chiar.mo Prof." / "Chiar.ma Prof.ssa"
%    * oppure "Prof." / "Prof.ssa" / "Dott." / "Dott.ssa"
%   a seconda di quanto concordato con il/la relatore/relatrice.
% - "NOME RELATORE" con Nome e Cognome del relatore.
{\large{\bf Relatore:\\
Chiar.ma Prof.\\
NOME RELATORE}} \\[5mm]

% ------------------------------------------------------------------
% CORRELATORE (FACOLTATIVO)
% ------------------------------------------------------------------
% Se HAI un correlatore:
%   - Sostituisci "Chiar.mo Prof." con il titolo accademico corretto.
%   - Sostituisci "NOME CORRELATORE" con Nome e Cognome del correlatore.
%
% Se NON hai un correlatore:
%   - Commenta il blocco qui sotto anteponendo % a ciascuna riga,
%     oppure cancellalo.
\large{\bf Correlatore: \\
Chiar.mo Prof. \\
NOME CORRELATORE}
\end{minipage}
\hfill
\begin{minipage}[t]{0.47\textwidth}\raggedleft
% ------------------------------------------------------------------
% NOME DELLAUREANDO / LAUREANDA
% ------------------------------------------------------------------
% Sostituire "NOME LAUREANDO" con:
%   Nome Cognome
% esattamente come compare nei documenti ufficiali.
{\large{\bf Presentata da:\\
NOME LAUREANDO}}
\end{minipage}

\vspace{10mm}
\begin{center}
% ------------------------------------------------------------------
% APPELLO, DATA E ANNO ACCADEMICO
% ------------------------------------------------------------------
% "NUMERO Appello":
%   - Sostituire con il numero dell'appello di laurea (es. "I Appello", "II Appello")
%     secondo le indicazioni della Segreteria/CdS.
%
% "Sessione MESE ANNO":
%   - Sostituire con la data ufficiale della seduta di laurea
%
% "Anno Accademico 20xx/20xx":
%   - Sostituire con l'Anno Accademico corretto (es. "Anno Accademico 2025/2026").
{\large{\bf Sessione MESE ANNO \\}}
{\large{\bf Anno Accademico 20xx/20xx\\}}
\end{center}
\end{titlepage}

\end{document}
